\documentclass[11pt,a4paper]{article}
\usepackage[margin=2.5cm]{geometry}
\usepackage{parskip}
\pagestyle{empty}
\begin{document}
\begin{center}
{\large\bfseries Within-Family Gains vs External Skill in Fractional SEIT Tuberculosis Forecasts}\\[0.6em]
Amaury de Souza and Francisco Federico Garc\'{\i}a Crespi
\end{center}
\medskip
\noindent\textbf{Abstract}\par
An improvement over an integer-order epidemic model does not necessarily imply useful forecasts. We investigated this distinction by independently reimplementing a fractional-order susceptible--exposed--infectious--treated (SEIT) tuberculosis model using 264 monthly notifications from Brazil (2001--2022). Its integer-order counterpart was refitted independently under the same calibration conditions. We then compared both models with persistence, seasonal naive, and seasonal autoregressive integrated moving average (SARIMA) forecasts across 13 expanding-window origins and leads of 1--12 months. Fractional SEIT achieved lower observed root mean squared error (RMSE) and mean absolute error than integer SEIT at every lead (RMSE at lead 1: 792 versus 1186 cases per month). Yet persistence (390) and SARIMA (536) had lower errors throughout, and the fractional model improved on seasonal naive (1091) only at leads 1--5. The fractional order and the progression rate lay at their lower search bounds; in a single-run sensitivity analysis, lowering the bounds moved both estimates and narrowed the within-family margin without changing its sign. Transmission, removal, and excess-mortality rates varied substantially among near-equivalent fits, although their predicted trajectories were similar. The basic reproduction number and disease-free equilibrium were evaluated only for an autonomous reference, not as properties of the time-varying forecasting system. For this national series, with dependent forecast errors and no formal significance tests, the within-family improvement did not establish competitive forecasting performance against persistence or SARIMA.

\medskip
\noindent\textbf{Keywords:} Mathematical epidemiology; Fractional-order differential equations; Practical identifiability; Forecast evaluation; Basic reproduction number; Tuberculosis
\end{document}
